% !PRINTED 203
\weberpage{203}
%<<CONT
can be determined, then $\alpha^2$, for example, can be computed from the following quotient:
\[
\frac{\alpha\beta\,.\,\alpha\gamma}{\beta\gamma}
\]

It now remains, therefore, to form the cubic equation whose
roots are $a_2', a_2'', a_2'''$. But this is obtained at once from
the invariants.

If one forms the invariants $A', B'$ from the transformed
form (4), then by §.\,64, (8), (9), (12), since $r = 1$, it follows that
\weq{11}{\begin{gathered}
A = a_2'^{\,2} - 3\,a_1'\,a_3',\\
B = 2\,a_2'^{\,3} - 9\,a_1'\,a_2'\,a_3',\\
D = a_1'^{\,2}\,a_2'^{\,2}\,a_3'^{\,2} - 4\,a_1'^{\,3}\,a_3'^{\,3},
\end{gathered}}
and from these, by elimination of $a_1'\,a_3'$,
\[
a_2'^{\,3} - 3\,a_2'\,A + B = 0.
\]
Hence $a_2', a_2'', a_2'''$ are the roots of the equation, cubic in $z$,
\weq{12}{z^3 - 3\,A\,z + B = 0.}

This is surely the simplest form that can be given to the cubic
resolvent of the biquadratic equation.

\begin{center}
§. 66.

\gesperrt{The Covariants.}
\end{center}

If we form the covariant $H$, according to §.\,64, (2), for the transformed
 form
\weq{1}{f(x,\,y) = F(\xi,\,\eta) = a_1'\,\xi^3\eta + a_2'\,\xi^2\eta^2 + a_3'\,\xi\eta^3}
we obtain
\weq{2}{\begin{gathered}
-H = 3\,a_1'^{\,2}\,\xi^4 + 4\,a_1'\,a_2'\,\xi^3\eta - (6\,a_1'\,a_3' - 4\,a_2'^{\,2})\,\xi^2\eta^2\\
{} + 4\,a_2'\,a_3'\,\xi\eta^3 + 3\,a_3'^{\,2}\,\eta^4,
\end{gathered}}
and from this it follows
\weq{3}{H + 4\,a_2'\,f = -3\,(a_1'\,\xi^2 - a_3'\,\eta^2)^2.}

This combination is thus, apart from the factor $-3$, the
square of a quadratic form, and the same holds if we replace
$a_2'$ by $a_2'', a_2'''$. We therefore set, as an
abbreviation,
%CONT>>
